\section{Action ellipse, radian, and double holonomy return}
\label{sec:elipse-accion-holonomia-rev8}
\index{action!ellipse}
\index{radian!action sector}
\index{Möbius!double return}

The preceding focal theorem fixes the angular form and the determinantal reader
fixes the action scale. Let us define
\[
 \eta=\operatorname{artanh}\frac{C^*}{A},
 \qquad
 a_S=\sqrt{2\hbar_{\rm HMT}^{\rm ret}e^\eta},
 \qquad
 b_S=\sqrt{2\hbar_{\rm HMT}^{\rm ret}e^{-\eta}}.
\]
The symplectic orbit
\[
 Q(\theta)=a_S\cos\theta,
 \qquad
 P(\theta)=b_S\sin\theta
\]
simultaneously preserves the shape ratio of the angular block and the area
fixed by the action line.

\begin{theorem}[Phase swept area]
\label{thm:area-fase-elipse-rev8}
For every \(\theta\),
\[
 \operatorname{Area}(\theta)
 =\frac12\int_0^\theta(QP'-PQ')\,dt
 =\hbar_{\rm HMT}^{\rm ret}\theta.
\]
Therefore,
\[
 \operatorname{Area}(1\ {\rm rad})=\hbar_{\rm HMT}^{\rm ret},
 \qquad
 \operatorname{Area}(2\pi)=h_{\rm HMT}^{\rm ret}.
\]
\end{theorem}

\begin{proof}
Since \(QP'-PQ'=a_Sb_S\) and
\(a_Sb_S=2\hbar_{\rm HMT}^{\rm ret}\), the integral equals
\(\hbar_{\rm HMT}^{\rm ret}\theta\). The second identity uses
\(h_{\rm HMT}^{\rm ret}=2\pi\hbar_{\rm HMT}^{\rm ret}\).
\end{proof}

APP supplies the arean oriented; \(A,C^*\) fix the form and the foci; the branch
determinantal fixes the scale. The turn primitive encloses \(h\) and a radian
of phase sweeps \(\hbar\). The radian ordinary remains
\(1\,\mathrm{rad}=180^\circ/\pi\); the construction adds its genealogy as
sector unitary of action, not another constant angular.

For a persistent section with holonomy
\(\varepsilon_\gamma\in\{+1,-1\}\), the condition
\[
 \varepsilon_\gamma e^{iJ_\gamma/\hbar}=1
\]
distinguishes two closures. In the cylinder, \(\varepsilon_\gamma=+1\), and the primitive
cycle transports \(J=2\pi\hbar=h\). In the Möbius band,
\(\varepsilon_\gamma=-1\): one visible turn changes sheet and transports
\(J=\pi\hbar=h/2\); the second restores orientation and completes \(h\). Per
radian of the internal phase, \(\hbar\) is transported in both cases; the factor
\(1/2\) belongs to the projection onto the visible turn of the Möbius base.

\begin{figure}[htbp]
\centering
\resizebox{.98\textwidth}{!}{\begin{tikzpicture}[
  font=\scriptsize,
  box/.style={draw,rounded corners=2mm,align=center,minimum height=.9cm,
    inner xsep=5pt},
  arr/.style={-{Stealth[length=2.2mm]},thick}]

  \begin{scope}[shift={(0,0)}]
    \draw[->,hmtgray] (-3.0,0)--(3.0,0) node[right]{\(u_+\)};
    \draw[->,hmtgray] (0,-2.0)--(0,2.0) node[above]{\(u_-\)};
    \draw[hmtblue,very thick] (0,0) ellipse (2.65cm and 1.45cm);
    \fill[hmtgold] (-1.85,0) circle (2pt) node[below]{\(F_-^{(\lambda)}\)};
    \fill[hmtgold] (1.85,0) circle (2pt) node[below]{\(F_+^{(\lambda)}\)};
    \node[text=hmtblue,font=\bfseries] at (0,2.55) {yield ellipse};
    \node[align=center,text=hmtgray] at (0,-2.4)
      {\(q_+=e^{-a_\lambda^2}\), \(q_-=e^{-b_\lambda^2}\)};
  \end{scope}

  \draw[arr,hmtgold] (3.3,0)--node[above,align=center]{same form\\action scale} (6.0,0);

  \begin{scope}[shift={(9.0,0)}]
    \draw[->,hmtgray] (-2.8,0)--(2.8,0) node[right]{\(Q\)};
    \draw[->,hmtgray] (0,-2.0)--(0,2.0) node[above]{\(P\)};
    \fill[hmtlightgold] (0,0)--(2.5,0)
      arc[start angle=0,end angle=57.2958,x radius=2.5cm,y radius=1.35cm]--cycle;
    \draw[hmtviolet,very thick] (0,0) ellipse (2.5cm and 1.35cm);
    \draw[hmtgold,thick] (0,0)--(2.5,0);
    \draw[hmtgold,thick] (0,0)--({2.5*cos(57.2958)},{1.35*sin(57.2958)});
    \node[text=hmtviolet,font=\bfseries] at (0,2.55) {action ellipse};
    \node[align=center] at (.7,.48) {\(\theta=1\)\\area \(\hbar\)};
    \node[align=center,text=hmtgray] at (0,-2.4) {vuelta \(2\pi\): area \(h\)};
  \end{scope}

  \node[box,draw=hmtblue,fill=hmtlightblue,minimum width=3.2cm] (cyl) at (3.2,-4.3)
    {cilindro \(\varepsilon=+1\)\\\(2\pi\mapsto h\)};
  \node[box,draw=hmtviolet,fill=hmtlightviolet,minimum width=4.2cm] (mob) at (9.0,-4.3)
    {Möbius case \(\varepsilon=-1\)\\\(2\pi\mapsto h/2\), \(4\pi\mapsto h\)};
  \draw[arr,hmtgray] (cyl)--node[
    above,align=center,text width=1.9cm,fill=white,inner sep=1pt
  ]{change of\\holonomy} (mob);
  \node[align=center,text=hmtgray,font=\footnotesize] at (6.1,-5.55)
    {The internal phase transports \(\hbar\) per radian;\\the factor two belongs to the visible turn.};
\end{tikzpicture}
}
\caption{The ellipse of rates fixes the form; its lift to the action line fixes
the area. The cylindrical closure and the double Möbius return distinguish the internal phase
from the visible turn.}
\label{fig:elipse-accion-holonomia-rev8}
\end{figure}

\subsection{Hilbert--Beltrami--Klein Geometry on the Action Ellipse}
\label{sec:hilbert-beltrami-klein-elipse-accion}

The same ellipse that carries the symplectic area admits an intrinsic projective
geometry. Let
\[
 \mathbb E_{\rm act}
 :=\left\{(Q,P)\in\mathbb R^2:
 \frac{Q^2}{a_S^2}+\frac{P^2}{b_S^2}<1\right\},
 \qquad
 \mathcal L_{\rm act}(Q,P)
 :=\left(\frac{Q}{a_S},\frac{P}{b_S}\right).
\]
The linear map \(\mathcal L_{\rm act}\) identifies
\(\mathbb E_{\rm act}\) with the unit disk \(\mathbb D\). For two
distinct points \(x,y\in\mathbb E_{\rm act}\), let \(a,b\in\partial
\mathbb E_{\rm act}\) be the intersections of the line \(xy\) with the boundary,
ordered as \(a,x,y,b\). We define
\begin{equation}
 d_{\rm H}^{\rm act}(x,y)
 :=\frac12\log
 \frac{\lVert y-a\rVert\,\lVert x-b\rVert}
      {\lVert x-a\rVert\,\lVert y-b\rVert}.
 \label{eq:distancia-hilbert-elipse-accion}
\end{equation}
The expression is the cross-ratio of the four collinear points and is therefore
independent of the auxiliary norm used on that line.

\begin{proposition}[Beltrami--Klein realization of the action ellipse]
\label{prop:beltrami-klein-elipse-accion}
The pair \((\mathbb E_{\rm act},d_{\rm H}^{\rm act})\) is isometric, through
\(\mathcal L_{\rm act}\), to the Beltrami--Klein model of the hyperbolic plane.
Its unoriented geodesics are exactly the open chords of the
ellipse. The boundary \(\partial\mathbb E_{\rm act}\) simultaneously preserves, in a
typed manner, the symplectic orbit
\[
 \gamma_{\rm act}(\theta)
 =(a_S\cos\theta,b_S\sin\theta),
 \qquad
 \frac12\int_0^\theta(QP'-PQ')\,dt
 =\hbar_{\rm HMT}^{\rm ret}\theta.
\]
\end{proposition}

\begin{proof}
Projective transformations preserve cross-ratios. Since
\(\mathcal L_{\rm act}\) is an invertible linear transformation that carries
the ellipse to the unit disk, it transports
\eqref{eq:distancia-hilbert-elipse-accion} to the Hilbert distance of the
disk. This is the Beltrami--Klein metric, whose geodesics are its chords.
The symplectic identity of the boundary is
\cref{thm:area-fase-elipse-rev8} written in the normalized chart.
\end{proof}

The complete chart is thereby typed by
\[
 \mathfrak R_{\rm act}:
 (A,C^*,\hbar_{\rm HMT}^{\rm ret})
 \longmapsto
 (\mathbb E_{\rm act},d_{\rm H}^{\rm act},
  dP\wedge dQ,\gamma_{\rm act}).
\]
The projective metric governs the interior; the symplectic form and the area
law govern the traversed boundary. A realization as a physical region or
as a domain of covariance matrices is formulated through a subsequent functor
that explicitly preserves these two structures; the preceding construction
already fixes its mathematical domain and the invariants that this functor must preserve.
